\section{Induction for complete intersections}
\label{ind:section}

We apply the transport theorem with all the tests allowed in
\cref{sec:conventions}. Throughout this section Virasoro means the
ordinary descendant constraints, with the first Hodge grading and the
super conventions fixed there. The two results used from the preceding
sections are the tested transport statement \cref{sew:transport} and the
primitive detection statement \cref{det:detection}.

\subsection{The toric bundle input with the first Hodge grading}

We first specify the form of the standard toric bundle result needed
below. Let $B$ be a smooth projective variety, let $L_1,\ldots,L_N$ be line
bundles on $B$, and let $E\to B$ be a smooth projective toric bundle
obtained by taking, fiberwise, the toric quotient of
$L_1\oplus\cdots\oplus L_N$ by a fixed subtorus. The toric fiber is
assumed smooth and projective. These are the bundles of
\cite[\S1.6]{CGT}; in particular they include projectivizations of sums
of line bundles. Write $s$ for the number of torus fixed points in the
fiber. The fixed locus in $E$ consists of $s$ sections isomorphic to
$B$.

Theorem~1.4 of \cite{CGT} constructs a symplectic loop transformation
$M$, with a nonequivariant limit, for which
\begin{equation}
\label{ind:descendant-transform}
 Z_E=c\,\widehat M\, Z_B^{\otimes s}.
\end{equation}
Here $c$ is a nonzero scalar independent of the descendant variables;
its value is immaterial to the constraints. This is an identity of total
descendant potentials on full super cohomology. The classical
transformation respects the grading operators. We explain the passage
from the half-total-degree notation in that paper to our first-index
convention, including the normalization after quantization.

\begin{proposition}[Toric bundle transfer]
\label{ind:toric}
Suppose that the first-Hodge Virasoro constraints hold for $B$, in all
genera and with all descendant insertions. They then hold for $E$ with
the same scope. The assertion can be iterated for towers of the toric
bundles described above.
\end{proposition}

\begin{proof}
For a smooth projective target $Y$, set
\[
 \mu^{\mathrm{tot}}_Y\big|_{H^{p,q}(Y)}
       =\frac{p+q-\dim Y}{2},\qquad
 \nu_Y\big|_{H^{p,q}(Y)}=\frac{p-q}{2}.
\]
Thus $\mu_Y=\mu_Y^{\mathrm{tot}}+\nu_Y$. Divisors, line-bundle Chern
classes, and toric and equivariant parameters have Hodge difference
zero. A Gromov--Witten correspondence preserves total Hodge difference:
its algebraic virtual class, descendant powers, and evaluation
pushforward have diagonal Hodge type. Consequently the transformations
built from these correspondences and characteristic classes are
equivariant for Hodge difference.

There is a small parameter issue in applying this observation.
Fundamental solutions at an arbitrary auxiliary parameter are
Hodge-equivariant as \emph{families}, with the corresponding action on
the parameter. We need an intertwining identity on cohomology at a
fixed parameter. Choose the auxiliary base parameter $\tau_B=0$ and
retain only the toric divisor parameters. Such a specialization is
permitted: the total descendant potentials are independent of the
auxiliary parameters, and a single auxiliary value suffices in the
construction \cite[\S4, footnote~11]{CGT}. It imposes no restriction on
the descendant variables in \eqref{ind:descendant-transform}.

For completeness, this equivariance can be checked on the final
nonequivariant transformation rather than on separate singular
localization factors. In the notation of \cite[\S4.4]{CGT}, the
transformation is the product of the matrix $T$ with the inverse of the
stationary-phase matrix of the oscillating integral $\mathcal I$;
write $\mathcal I_{\mathrm{as}}$ for that matrix. Both matrices have
columns indexed by a toric basis class times a homogeneous base class
$\phi_b$. On this common column space define $\nu_{\mathrm{col}}$ to
have eigenvalue $(p_b-q_b)/2$ when $\phi_b$ has type $(p_b,q_b)$.
The toric basis classes have Hodge difference zero.

At $\tau_B=0$, every column of the base fundamental solution has the
Hodge difference of its input, because its coefficients are algebraic
Gromov--Witten correspondences. The toric operations producing $T$
preserve this difference. The construction of $\mathcal I_{\mathrm{as}}$
uses the same base solution and differentiates it only in the
line-bundle Chern-class directions. These directions, the toric
differential operators, and the scalar stationary-phase coefficients
all have Hodge difference zero. Consequently
\[
 \nu_E T=T\nu_{\mathrm{col}},\qquad
 \nu_{B^{\sqcup s}}\mathcal I_{\mathrm{as}}
       =\mathcal I_{\mathrm{as}}\nu_{\mathrm{col}}.
\]
These identities include the columns indexed by odd classes. Forming
$T\mathcal I_{\mathrm{as}}^{-1}$ and taking its nonequivariant limit
$M$ gives
\begin{equation}
\label{ind:nu-intertwining}
 \nu_E M=M\nu_{B^{\sqcup s}}.
\end{equation}
The classical grading identity proved in
\cite[\S4.4, Equations~(42) and~(44)]{CGT} is
\[
 \left(z\partial_z+\frac12+\mu_E^{\mathrm{tot}}
          +\frac{c_1(E)\cup}{z}\right)M
 =M\left(z\partial_z+\frac12+\mu_{B^{\sqcup s}}^{\mathrm{tot}}
          +\frac{c_1(B^{\sqcup s})\cup}{z}\right).
\]
Adding \eqref{ind:nu-intertwining} gives the same identity with $\mu$ in
place of $\mu^{\mathrm{tot}}$. Since $M$ commutes with multiplication by
$z$, it intertwines all the classical positive operators obtained from
$l_0(zl_0)^k$.

We use the central normalization for this new grading. For a target
$Y$ of dimension $d$, the scalar in $L_0$ is
\begin{equation}
\label{ind:central-normalization}
 C_Y=\frac{\chi(Y)}{16}-\frac14\str(\mu_Y^2)
 =\frac1{48}\int_Y\bigl((3-d)c_d(Y)
                  -2c_1(Y)c_{d-1}(Y)\bigr).
\end{equation}
The equality is the Hodge-index Riemann--Roch identity, and the
first-Hodge grading equation $L_0Z_Y=0$ is the usual Hori equation
\cite[Theorem~2.1 and Proposition~2.6]{Getzler1999}. For a
zero-dimensional target the second Chern-number term is omitted.
No half-total-degree grading equation is being asserted here.

Quantization is in the super symplectic space; the resulting cocycle
uses supertrace. Conjugation of the first-Hodge quantized operators by
$\widehat M$ can differ from the target operators only by scalars.
For $k=0$ the scalar difference is zero, because both operators
annihilate the same nonzero potential in
\eqref{ind:descendant-transform}, by their first-Hodge grading
equations. For $k\ne0$ the commutator with $L_0$ forces the scalar
difference to vanish. This is precisely the argument of
\cite[Proposition~1.3]{CGT}, now applied with the first-Hodge operators
and \eqref{ind:central-normalization}. In particular one need not, and
in general cannot, identify $C_E$ with $sC_B$: the quantization cocycle
accounts for their difference. Equation
\eqref{ind:descendant-transform} now transfers all constraints. The
same argument applies at each stage of a tower.
\end{proof}

\subsection{The splitting degeneration}

Let $X\subset\PP^{n+r}$ be a smooth complete intersection with positive
degrees $(d_1,\ldots,d_r)$, and choose $d_r=a+b$ with $a,b>0$.
After a smooth deformation we can take general defining equations
$f_1,\ldots,f_r$ and general forms $g_a,g_b$ of degrees $a,b$.
Consider the family
\begin{equation}
\label{ind:pencil}
 f_1=\cdots=f_{r-1}=0,\qquad g_a g_b=t f_r.
\end{equation}
Its two central components before resolution are
\[
 U=\{f_1=\cdots=f_{r-1}=g_a=0\},\qquad
 M=\{f_1=\cdots=f_{r-1}=g_b=0\}.
\]
Their intersection is the smooth complete intersection
\[
 D=\{f_1=\cdots=f_{r-1}=g_a=g_b=0\}.
\]
The singular locus of the total family is
\[
 S=\{f_1=\cdots=f_{r-1}=g_a=g_b=f_r=0\}.
\]
It has dimension $n-2$ when nonempty. Bertini gives the required
smoothness and transversality for these choices. In $M$ its ideal is
the regular sequence $(g_a,f_r)$, so
\begin{equation}
\label{ind:split-normal}
 N_{S/M}\simeq\cO_S(a)\oplus\cO_S(d_r).
\end{equation}

Transverse to $S$ the local equation of \eqref{ind:pencil} is
$xy=tz$. Blowing up the ideal of the appropriate central component
resolves this ordinary double-point family. The resulting total space
is smooth, its central fiber is a transverse union
\begin{equation}
\label{ind:resolved-splitting}
 U\cup_D\widetilde M,\qquad \widetilde M=\Bl_S M,
\end{equation}
and the strict transform of $D$ is again $D$, since $S$ is a Cartier
divisor in $D$. The two normal bundles of the seam are dual. The
construction is projective, being a blowup of a projective family.
These are also the splitting data in
\cite[Theorem~5.3 and its proof]{ABPZ}; we use that result here for the
geometry and its curve-class qualifications, not as a Virasoro theorem.
If $S$ is empty, no blowup is necessary. In dimension one $D$ can be a
disconnected zero-dimensional smooth scheme and $S$ is empty.

The two ends $U$ and $M$ have smaller total degree than $X$. The seam
$D$ and center $S$ have smaller dimension. Thus this construction is
compatible with induction first on dimension and then, at fixed
dimension, on
\begin{equation}
\label{ind:complexity}
 \kappa(d_1,\ldots,d_r)=\sum_i(d_i-1).
\end{equation}
Linear equations can be eliminated and do not change $\kappa$.
Replacing $d_r$ by either $a$ or $b$ strictly decreases $\kappa$.

\subsection{Coherent classes for the blowup step}

For a positive-dimensional connected complete intersection $Y$, let
$A_Y\subset H^*(Y,\Q)$ denote the image of projective-space cohomology.
For a nonempty zero-dimensional complete intersection, let $A_Y$ be
the span of its unit, that is, the sum of the units of its points.
These subspaces are nondegenerate for their induced Poincare pairings.
Their complements occur only in middle cohomology, by weak and hard
Lefschetz. The zero-dimensional convention will be useful when several
points form a single center.

\begin{lemma}[Blowup detection data]
\label{ind:blowup-data}
Let $M$ be an $n$-dimensional smooth complete intersection and let
$S\subset M$ be a smooth codimension-two complete intersection cut by
two ambient equations. Write $p:\widetilde M=\Bl_S M\to M$,
$j:E\hookrightarrow\widetilde M$, and $\pi:E\to S$. Then
\begin{equation}
\label{ind:blowup-B0}
 B_0=p^*A_M\oplus j_*\pi^*A_S
\end{equation}
is a nondegenerate rational Hodge subspace generated by diagonal Hodge
classes. It contains all cohomology outside $H^n(\widetilde M)$ and
all nonprimitive middle classes for a polarization
$p^*H-\epsilon E$, with $\epsilon>0$ sufficiently small and rational.
Its indicated generators have coherent single-class lifts in the
blowup configuration of \cref{sew:transport}.
\end{lemma}

\begin{proof}
The codimension-two blowup formula gives
\begin{equation}
\label{ind:blowup-cohomology}
 H^k(\widetilde M,\Q)
  =p^*H^k(M,\Q)\oplus j_*\pi^*H^{k-2}(S,\Q).
\end{equation}
The second summand has Hodge shift $(1,1)$. The only possible
nonambient summands in \eqref{ind:blowup-cohomology} therefore occur
for $k=n$. Put $\xi=c_1(\cO_E(1))$, with
$N_{E/\widetilde M}=\cO_E(-1)$. The push-pull and self-intersection
formulas give, in complementary degrees,
\begin{align}
 \int_{\widetilde M}p^*u\,j_*\pi^*a&=0,\label{ind:orthogonal}\\
 \int_{\widetilde M}j_*\pi^*a\,j_*\pi^*b
       &=-\int_E\xi\,\pi^*(ab)=-\int_Sab.
       \label{ind:exceptional-pairing}
\end{align}
For the first equality, the integrand on $E$ is pulled back from $S$
and has no fiber hyperplane factor. These identities prove
nondegeneracy of \eqref{ind:blowup-B0}.

The subspace $B_0$ is stable under multiplication by the ambient
hyperplane and by $E$. Indeed the projective bundle relation expresses
powers of $\xi$ through Chern classes of $N_{S/M}$, and those Chern
classes are ambient; the pushforward of $S$ in $M$ is ambient as well.
Equivalently the blowup ring formula closes on the two ambient
summands in \eqref{ind:blowup-B0}. For small positive rational
$\epsilon$, $p^*H-\epsilon E$ is ample. Since $H^{n-2}(\widetilde M)$
is contained in $B_0$, so is its image under this Lefschetz operator.
That image is the nonprimitive part of $H^n(\widetilde M)$.
Orthogonality then implies that $J=B_0^\perp$ is primitive middle
cohomology, with its Hodge--Riemann polarization.

When $n=2$ and $S=\{s_1,\ldots,s_m\}$, the exceptional part of
$B_0$ is generated by $E_1+\cdots+E_m$. Formula
\eqref{ind:exceptional-pairing} has matrix $-I_m$ on the individual
exceptional classes. The omitted differences, with coefficient sum
zero, are primitive for $p^*H-\epsilon\sum_iE_i$. They are included
in $J$, not discarded. This also proves directly that the ambient
constant on a zero-dimensional center suffices in
\eqref{ind:blowup-B0}.

It remains to describe the lifts needed for transport. Pullbacks of
ambient classes on $M$ extend by pullback through the normal-cone
families and their inserted ruled levels. Each exceptional generator
in \eqref{ind:blowup-B0} is $E$ multiplied by an ambient pullback.
In the normal-cone expansion along $E$, the divisor $E$ follows into
the nonjoining section of the final ruled component. Use the class
Gysin-pushed from that section, multiplied by the corresponding
ambient class. It is disjoint from the joining boundary, and hence its
restriction there is zero. Give it zero restriction on all other
pieces. These are the specializations of the indicated absolute
classes, as can be seen from the strict transform of
$E\times\A^1$ in the deformation to its normal cone. The same
construction after inserting additional levels places the class at
the unchanged final section. At a surgery it remains on the segment
containing that end, and is zero on segments from which the end was
removed. Thus all retained component restrictions agree in the
comparisons required by \cref{sew:transport}.
\end{proof}

\subsection{Curve classes and deformation}

The degree statement in the transport theorem concerns the divisors
which extend through the families. We record why sufficiently many
such divisors are available here.

\begin{lemma}[Degree separation]
\label{ind:degrees}
The ordinary and blowup configurations used in
\eqref{ind:resolved-splitting} can be equipped with coherent numerical
divisor tests which distinguish every curve class with a possibly
nonzero ordinary Gromov--Witten contribution, after choosing very
general smooth fibers when necessary. The same conclusion holds for
disconnected total classes.
\end{lemma}

\begin{proof}
If a positive-dimensional complete intersection has dimension different
from two, its integral curve homology is detected by the hyperplane
degree. For curves this follows from $H_2(Y,\Z)=\Z$; for dimension
at least three it follows from the integral Lefschetz theorem.
There is no torsion ambiguity. Complete-intersection surfaces are
simply connected, so Poincare duality and the universal coefficient
theorem also show that their $H_2$ is torsion free.

For a complete-intersection surface with all linear equations
eliminated, the very general rational algebraic cohomology is ambient
unless the surface is $\PP^2$, a quadric, a cubic, or an intersection
of two quadrics. This is the complete-intersection Noether--Lefschetz
statement used in \cite[\S5.2]{ABPZ}. In the nonexceptional case,
ambient degree therefore distinguishes all effective classes on a
very general fiber. A class with a nonzero primitive part is not
algebraic there and cannot carry a stable map. Deformation invariance
makes its ordinary invariant zero also on any special smooth fiber,
with the curve class and insertions transported locally. An effective
disconnected total is itself an algebraic class, so this argument
applies to disconnected coefficients as well.

For the exceptional surfaces, all of $H^2$ has type $(1,1)$. Their
primitive integral lattice is negative definite by Hodge--Riemann.
The monodromy, which preserves that lattice and the hyperplane class,
is consequently finite. The local monodromy of
\eqref{ind:resolved-splitting} is also unipotent, since the total space
is smooth and the central fiber is semistable. It is therefore the
identity. In these exceptional-surface splittings the seam $D$ is a
line or a smooth conic, hence $H^1(D,\Q)=0$; the projective-plane base
case needs no splitting. The Mayer--Vietoris sequence consequently
identifies $H^2(U\cup_D\widetilde M,\Q)$ with the kernel of the
difference of the two restrictions to $H^2(D,\Q)$. The invariant cycle
theorem, with trivial monodromy, supplies a central-fiber lift of every
rational divisor class on the nearby surface. Its image in this
kernel is a matching pair
\begin{equation}
\label{ind:matching-divisors}
 (\delta_U,\delta_{\widetilde M})\in
 H^2(U,\Q)\oplus H^2(\widetilde M,\Q),\qquad
 \delta_U|_D=\delta_{\widetilde M}|_D,
\end{equation}
modulo the opposite boundary-divisor relation. One way to see the
Hodge type of these representatives is to use the column-zero
weight-complex description: restriction and Gysin are Hodge
morphisms, and strictness permits a type-$(1,1)$ representative for a
pure invariant divisor class. Equivalently, one may take closures of
invariant divisors in the resolved family.

These representatives extend coherently across surgery. Keep the
class on each retained end, and on every ruled neck and cap take the
pullback of the common class on $D$. On each shortened smoothing these
are the usual extensions in the deformation to a normal cone. Opposite
boundary twists change a representative, but their contributions to
the sum of degrees cancel at a matching contact. A basis of
$H^2(X,\Q)$ among these tests distinguishes all integral curve
classes on the exceptional surface, since the latter lattice has no
torsion.

Finally, for $\widetilde M=\Bl_S M$, the integral curve lattice is
\[
 H_2(\widetilde M,\Z)
       \simeq H_2(M,\Z)\oplus H_0(S,\Z).
\]
Retain the needed divisor degrees from $M$ by pullback and the degrees
against each connected component of the exceptional divisor. The
latter are represented by the nonjoining sections at the final end
of the blowup configuration; hence they are not altered by the cap
bookkeeping at an internal surgery. The pullback classes distinguish
the first summand, and exceptional intersection distinguishes the
second, with nonzero diagonal coefficient on its fiber generators.
If $M$ is a nonexceptional surface we choose it very general; if it is
exceptional we retain all its divisor degrees. The hyperplane degree
and these additional tests are among the compatible degrees of
\cref{sew:transport}. Finiteness at a fixed hyperplane and exceptional
bound is part of that theorem. This proves the claim.
\end{proof}

\begin{lemma}[Passage to special smooth fibers]
\label{ind:deformation}
Suppose the tested constraints, and hence the constraints detected by
\cref{det:detection}, hold on very general members of one of the
smooth complete-intersection families or nested blowup families above.
They then hold on every smooth member, for every locally transported
curve class.
\end{lemma}

\begin{proof}
Work locally in the smooth parameter space, using a topological
trivialization of the cohomology and curve-class local systems. The
ordinary Gromov--Witten evaluation tensors are flat morphisms of
Hodge structures with fixed Tate shifts. The cup pairing and the
ambient classes are flat, as are the ambient and exceptional
subspaces in \eqref{ind:blowup-B0}. Thus their orthogonal projectors
are flat. Connected components of a zero-dimensional center may be
labelled locally; their possible global permutation is irrelevant.

Fix a coefficient error and its Hermitian-square test from
\cref{det:detection}. All contractions in that test have balanced
Hodge type. The only nonflat ingredients in its expression in a
flat cohomology trivialization are first-Hodge-index projectors and
functions of them. The derivative of such a projector is
first-index off-diagonal: differentiating its idempotence and its
constant eigenspace decomposition shows that its diagonal blocks
vanish. Replacing one projector by its derivative in the balanced
contraction changes the total first Hodge degree while leaving all
other tensor morphisms at their fixed type. The resulting scalar is
zero. The product rule therefore shows that the Hermitian-square
scalar is locally constant. This is the same Hodge-balance argument
used for numerical tests in the degeneration calculation.

The square vanishes on the very general fibers under consideration,
so it vanishes throughout the connected smooth parameter space.
Positivity in \cref{det:detection} proves vanishing of the error on
each fiber, for all its primitive inputs. The ambient inputs follow
from the same scalar argument without free primitive slots. If a
curve class is nonalgebraic on a very general fiber, all its ordinary
Gromov--Witten tensors are zero there; the preceding reasoning still
applies. Coefficient extraction for the disconnected potential is
legitimate at each fixed polarized degree and genus power by the
finiteness conventions of \cref{sew:transport}.
\end{proof}

The general choices required here can be made simultaneously. Choose
the defining forms, factors, and smoothing form in the open parameter
space where all the indicated intersections are smooth. A fixed
nonzero smooth fiber of \eqref{ind:pencil} ranges over general
complete intersections when its smoothing equation ranges freely.
The ends and nested centers likewise range over their smooth complete
intersection parameter spaces. Removing the relevant countable
Noether--Lefschetz loci thus allows the very general degree assertions
at every surface stage. Subsequent passage to arbitrary smooth
members is supplied by \cref{ind:deformation}.

\subsection{Completion of the induction}
\label{ind:complete}

\begin{proof}[Proof of \cref{thm:main}]
A zero-dimensional smooth complete intersection is a disjoint union
of points. The point constraints are the Witten--Kontsevich theorem
\cite{Witten1991,Kontsevich1992}. For a disjoint union, the potential
is the product and each Virasoro operator is the sum of the component
operators, so the assertion follows in dimension zero.

Assume the assertion in dimensions below $n$. In dimension $n$,
induct on \eqref{ind:complexity}. Complexity zero is projective space
after eliminating linear equations. This is a toric bundle over a
point, so \cref{ind:toric} proves the base case. At positive
complexity choose $d_r\ge2$ and a splitting $d_r=a+b$. Use
\eqref{ind:resolved-splitting}. Both $U$ and $M$ have smaller
complexity, and $D,S$ have smaller dimension. Their full constraints
are therefore available by induction.

We first establish the constraints for $\widetilde M$. If $S$ is
empty, this is the assertion for $M$ already known. Otherwise set
\[
 E=\PP_S(N_{S/M}),\qquad
 C_E=\PP_E(\cO_E\oplus\cO_E(-1)),\qquad
 Q_S=\PP_S(N_{S/M}\oplus\cO_S).
\]
In the blowup configuration of \cref{sew:transport}, the main
smoothing is $\widetilde M$, and the three shorter target types are
$M,Q_S,C_E$. By \eqref{ind:split-normal}, $E$ and $Q_S$ are
projectivizations of sums of line bundles over $S$, and $C_E$ is one
more such projectivization over $E$. The constraints for $S$ and
\cref{ind:toric} give the full constraints for $Q_S$ and $C_E$.
Thus all shorter targets required by transport are known.

Apply \cref{sew:transport} to obtain every tested positive constraint
on $\widetilde M$. Use the hyperplane degree pulled back from $M$
and retain the additional degrees in \cref{ind:degrees}, including
individual exceptional degrees when the center is disconnected.
The coherent subspace in \cref{ind:blowup-data} meets all the
hypotheses of \cref{det:detection}. Detection proves every positive
constraint for $\widetilde M$, with arbitrary insertions. The degree
tests distinguish its relevant curve classes, and
\cref{ind:deformation} includes special smooth choices of the nested
intersections. The string and grading equations, or the string
equation and Virasoro commutators with the normalization
\eqref{ind:central-normalization}, supply the nonpositive constraints.

We now use the ordinary configuration for
$U\cup_D\widetilde M$. Its shorter targets are $U$, $\widetilde M$,
and the ruled neck
\[
 C_D=\PP_D(\cO_D\oplus N_{D/U}).
\]
The first two have just been treated and $C_D$ has the constraints by
\cref{ind:toric} and the induction hypothesis for $D$. All these
statements include all descendants, so they remain valid after every
allowed test or auxiliary insertion used by transport.

Take $B_0=A_X$ on the smooth main target. Weak and hard Lefschetz show
that $B_0$ contains the nonmiddle cohomology and the nonprimitive
middle part, and that $B_0^\perp$ is primitive. Its classes are
coherent: restrict ambient projective classes to the two ends, and
pull their common restriction on $D$ to every neck and cap. These are
the natural extensions in the pencil and the normal-cone families.
Theorem~\ref{sew:transport}, followed by \cref{det:detection}, proves
the positive constraints for $X$ for every coefficient measured by
the compatible degrees. By \cref{ind:degrees} this gives the
constraints per curve class on the very general fibers; in the
exceptional surface cases all divisor degrees have been retained
throughout. Lemma~\ref{ind:deformation} then gives the conclusion
on every smooth complete intersection of the chosen multidegree.
The nonpositive constraints follow as in the blowup step.

Both induction parameters strictly decrease where invoked, and no
restriction on genus or on the number or parity of primitive
insertions was imposed. This proves the theorem.
\end{proof}
